\section{从单任务模型到 Foundation Models}

这一讲是 guest lecture，主讲人是 Ranjay Krishna。开场先做了一个很重要的对比：课程前半段学到的，多是为单个任务收集数据、训练单个模型、在固定测试集上评估的范式；而今天的重点，是把视觉系统推进到一个更通用的层次，也就是 foundation models。换句话说，问题不再是``每个任务都单独训练一个模型''，而是``能不能先训练一个大而通用的模型，再适配到很多任务''。\footnote{根据字幕 00:00:05--00:03:40，讲者把课程前半段的单任务流程和 foundation model 的预训练-迁移流程做了对比。}

\begin{figure}[H]
\centering
\includegraphics[width=0.95\textwidth]{slides-images/slide-002.jpg}
\caption{课程传统范式：每个任务都训练一个单独模型}
\end{figure}

\begin{figure}[H]
\centering
\includegraphics[width=0.95\textwidth]{slides-images/slide-003.jpg}
\caption{Foundation model 的核心思想：先在大规模多样数据上预训练，再迁移到多个下游任务}
\end{figure}

\subsection{Foundation model 到底是什么}

讲者没有给出一个教科书式定义，而是强调了一组更实用的特征。foundation model 不是一个严格的分类术语，而是一种系统设计取向：它应该足够通用、足够鲁棒、拥有较大的参数量、训练数据也足够大，而且通常采用某种自监督或弱监督目标。

\begin{importantbox}{Foundation model 的常见特征}
\begin{itemize}
    \item \textbf{通用性}：同一模型可适配很多不同任务。
    \item \textbf{规模}：参数量通常大，训练数据也大。
    \item \textbf{预训练优先}：先学表示，再做下游适配。
    \item \textbf{弱监督/自监督}：尽量减少对人工标签的依赖。
\end{itemize}
\end{importantbox}

这讲只讨论视觉和多模态系统，不展开语言模型本身，但语言模型的思路会被借用来构建视觉 foundation model，尤其是“如何把输入压成一组 token，再交给一个通用模型去自动完成”这条路线。

\begin{figure}[H]
\centering
\includegraphics[width=0.95\textwidth]{slides-images/slide-005.jpg}
\caption{foundation model 的谱系：语言、视觉、多模态、生成和 chaining}
\end{figure}

\subsection{为什么视觉也需要 foundation models}

视觉任务本身就很碎片化。分类、检测、分割、检索、问答、captioning，各自的数据集、损失函数和评估标准都不一样。如果继续用“一个任务一个模型”的思路，就会不断重复造轮子。foundation model 的价值，在于让同一套表示跨任务复用，而不是只在一个 benchmark 上刷分。\footnote{字幕 00:02:02--00:03:40 中，讲者明确说到它允许在很少数据甚至零数据的情况下适配新任务。}

\begin{figure}[H]
\centering
\includegraphics[width=0.95\textwidth]{slides-images/slide-021.jpg}
\caption{先预训练一个不依赖监督标签的编码器，再迁移到多个下游视觉任务}
\end{figure}

\begin{knowledgebox}{为什么这个方向重要}
\begin{itemize}
    \item 标签昂贵，而网络数据丰富。
    \item 下游任务不断变化，固定任务模型很难复用。
    \item 统一表示有利于检索、生成、分类和问答等任务的共享。
\end{itemize}
\end{knowledgebox}

\begin{figure}[H]
\centering
\includegraphics[width=0.95\textwidth]{slides-images/slide-001.jpg}
\caption{这门课的目标之一：把视觉任务从单点模型推进到可迁移的 foundation models}
\end{figure}

\subsection{从 self-supervised 到 multimodal}

讲者把本讲和前面的 self-supervised learning 联系起来：如果对比学习能把同一张图的不同增强拉近，那么把“图像”和“文本”也放进同一个对比目标里，就能把图文空间统一起来。这就是 CLIP 的出发点。\footnote{字幕 00:03:40--00:06:40 中，讲者先从 self-supervised 的 contrastive objective 过渡到图文对比学习。}

\begin{figure}[H]
\centering
\includegraphics[width=0.95\textwidth]{slides-images/slide-024.jpg}
\caption{CLIP 的基本结构：图像编码器和文本编码器分别输出向量，再用相似度做匹配}
\end{figure}

